% !PRINTED 411
\weberpage{411}
%<<CONT
tion $\Phi$, isomorphic (simply or multiply) with the alternating permutation
group of $n$ figures. Let the group of the linear substitutions
 $\chi(\vartheta)$ be denoted by $L$. Since it is
finite, it must be identical with one of the polyhedral groups
considered in the seventh chapter.

For simplification we now make use of the theorem proved
in \S.\,51, 2., according to which the group $L$ can be
so transformed that an arbitrary one of the non-identical
substitutions of $L$ becomes a multiplication. But a transformation
 of the group $L$ is equivalent to this, that for $\vartheta$
a linear function of $\vartheta$ is put, to which the same property
 belongs as to the original function $\vartheta$.

If, then, we denote by $\vartheta_\pi$ the function of the variables
$x_0, x_1, \ldots, x_{n-1}$ that arises from $\vartheta$ by application of the permutation $\pi$,
we can choose $\vartheta$ so that for a definite,
but arbitrary permutation $\pi$
\[
\vartheta_\pi = \varepsilon\vartheta
\]
holds. If one applies $\pi$ repeatedly, it follows that
\[
\vartheta_{\pi^2} = \varepsilon^2\vartheta,\ \vartheta_{\pi^3} = \varepsilon^3\vartheta, \ldots,
\]
and if therefore $p$ is the degree of the permutation $\pi$, then $\varepsilon$ is a
$p$th root of unity.

If first $n = 3$, the alternating group consists of
the powers of the cyclic permutation $\gamma = (0, 1, 2)$. We
may assume that the substitution corresponding to the permutation,
 $\chi(\vartheta)$, is multiplicative, that is, that $\vartheta_\gamma = \varepsilon\vartheta$, where $\varepsilon$
is a third root of unity. Hence $\vartheta^3$ is an alternating
(or symmetric) function. As resolvent with one parameter
 we thus obtain the pure equation
\[
\vartheta^3 = z,
\]
where $z$ is an alternating function. We may take for $\vartheta$,
say, the \gesperrt{Lagrange} resolvent $x_0 + \varepsilon x_1 + \varepsilon^2 x_2$, and
obtain the well-known reduction of the cubic equation to
a pure equation (Bd. I, \S.\,159).

We pass to the case $n = 4$, in which the alternating
 group $A$ contains the permutations
\[
1,\ \alpha_1 = (0, 1)(2, 3),\ \alpha_2 = (0, 2)(1, 3),\ \alpha_3 = (0, 3)(1, 2)
\]
which form a four-group $B$; besides these there occur
%CONT>>
